\chapter{Resultados}\label{ch:results}

Este capítulo apresenta as análises computadas e verificadas: as superfícies de viabilidade
do presente (\S\ref{sec:atlas-present}), o zoneamento agroambiental
(\S\ref{sec:zoning-results}), a mudança projetada pelo CMIP6 (\S\ref{sec:shift}) e a fase de
uso atual da terra e validação (\S\ref{sec:realized}). Todos os valores quantitativos provêm
das saídas verificadas do modelo. As tabelas de referência e suplementares e as figuras de
diagnóstico estão nos Anexos.

\section{Atlas de viabilidade multissegmentar do presente}\label{sec:atlas-present}

O mecanismo baseado em conhecimento (\S\ref{sec:suitability}) produziu, para cada um dos sete
segmentos, uma superfície de viabilidade do presente sobre GO~e~DF a 250~m. Cada superfície
traz uma classe da FAO (S1/S2/S3/N) e uma banda de robustez de sensibilidade dos pesos de
$\pm 20\%$. Três verificações confirmam o comportamento esperado do mecanismo: todas as faixas
de bandas situam-se dentro dos intervalos admissíveis; as máscaras de cobertura da terra
excluem, como esperado, as áreas urbanas de Goiânia e Brasília e os grandes reservatórios; e os
padrões espaciais conformam-se à expectativa agronômica, com a viabilidade de culturas
concentrando-se nos planaltos planos e de solo profundo e a viabilidade de conservação
crescendo no terreno íngreme, nativo e ripário. O resumo quantitativo mais compacto é a
viabilidade média de cada segmento por zona (Tabela~\ref{tab:zone-suit}, \S\ref{sec:zoning-results}),
acompanhada da parcela da terra avaliada de cada segmento que alcança a classe S2 da FAO ou
superior.

Soja e cana-de-açúcar pontuam mais alto nos planaltos mais frios, altos e argilosos (parcelas
S2+ de 62,2\% e 60,8\%) e caem acentuadamente em terreno íngreme ou arenoso. As outras culturas
anuais acompanham a confiabilidade da chuva e o calor acumulado, além do declive e da
fertilidade (S2+ 64,2\%). A pecuária é graduada pela produtividade de forragem, de modo que
pastagens genuinamente ruins caem para S3/N: sua parcela viável ou superior é de 63,1\%, mas
sua viabilidade média territorial fica \emph{abaixo} das culturas na terra boa.

A piscicultura tem 10,0\% do território em classe S2+ . Ancorar seu fator de proximidade
à água superficial \textbf{semipermanente} (mediana de 2,3~km), em vez da água permanente (mediana de
13,6~km), eleva a superfície acima do piso. É a água semipermanente, e não a permanente, que de
fato governa a escolha de sítio para aquicultura em pequenos reservatórios.

A superfície de valor de conservação é genuinamente informativa (S2+ 33,2\%). 
Ela agrega, de forma compensatória, a conectividade a áreas protegidas, a
heterogeneidade do terreno, o serviço de armazenamento de carbono, o serviço de água e a
marginalidade da terra. É bem graduada (média 0,45, desvio padrão espacial 0,12) e identifica
uma zona cujo melhor uso absoluto é a própria conservação --- o planalto alto, íngreme e rico
em carbono orgânico onde a viabilidade de culturas colapsa (\S\ref{sec:zoning-results},
Zona~6). A geração fotovoltaica atinge 27,3\% em S2+, discriminando as planícies mais baixas,
quentes e de alta irradiância dos planaltos mais frios. Essas planícies são favorecidas pela
irradiância e por um fator de índice de claridade, mas penalizadas por uma redução de
desempenho por calor no trimestre mais quente.

GO~e~DF é climaticamente quase uniforme, por isso as pertinências climáticas saturam, e a 
discriminação real entre terra boa e ruim vem do \textbf{solo e do terreno}. 
Uma decomposição da variância por fator torna essa leitura quantitativa 
(Figura~\ref{fig:factor-variance} do Anexo): decompondo a
variância espacial da viabilidade de cada segmento em parcelas por fator, os fatores climáticos
quase uniformes quase nada contribuem para a \emph{variação} entre terra boa e ruim. Para a
soja, o declive responde por 0,93 da variância e a argila por 0,07, com o clima $\approx$ 0. A
cana-de-açúcar se divide entre declive (0,59) e argila (0,41). A pecuária é governada pela
umidade do solo (0,58) e pelo carbono orgânico (0,25). E a conservação, pela conectividade a
reservas (0,43) e pela rugosidade do terreno (0,21), com a fração de vegetação nativa
respondendo por 0,16 e os fatores climáticos móveis em $\approx$ 0 no presente. Essa
concentração em um único fator para as culturas não decorre de peso AHP enviesado --- os pesos
de soja, cana-de-açúcar, outras culturas anuais e piscicultura são os \emph{mais equilibrados}
dentre os sete segmentos (razão máx./mín. de 1,4--1,7$\times$; Tabela~\ref{tab:spec-longtable}
do Anexo), inclusive mais que os da pecuária (4,25$\times$, deliberadamente reponderada para
fugir dos fatores climáticos saturados) --- mas da saturação espacial genuína já descrita acima.
Ressalva: para soja e outras culturas anuais, apenas argila (e declive) de fato discriminam
espaço --- pH e carbono orgânico do solo também saturam (parcela $\approx$ 0, apesar de não
serem climáticos), de modo que ``solo'' é, para essas duas culturas, um discriminador mais
estreito do que ``solo e terreno'' em conjunto. Os fatores climáticos atuam, assim, como
comportas de viabilidade: eles abrem ou fecham o intervalo possível, mas não desenham a
estrutura espacial fina do atlas, que repousa sobre solo e terreno --- uma dependência retomada
como limitação (\S\ref{sec:limitations}).

A decomposição por fator é agregada por pixel amostrado --- ela mostra \emph{quanto} cada fator
contribui em média, mas não \emph{onde} nem em qual escala espacial essa contribuição se
sustenta. Para responder a isso diretamente, o desvio-padrão local (janela móvel) de cada banda
padronizada da pilha de atributos foi calculado e agregado por tema --- clima
\emph{versus} terreno+solo --- em três raios de janela, incluindo 4,75~km, o tamanho da própria
célula nativa do TerraClimate (Figura~\ref{fig:scale-mismatch} e Tabela~\ref{tab:scale-mismatch}
do Anexo). Mesmo nessa escala --- a mais favorável possível ao clima, por ser a sua própria
resolução nativa --- a variabilidade espacial local de terreno+solo é 42$\times$ maior que a do
clima ($\log_2 = 5{,}40$); em janelas mais finas (750~m, 250~m) a razão sobe para 198$\times$ e
320$\times$. O mapa confirma que esse padrão é quase uniforme em todo o território, não um
artefato de algumas células: terreno e solo produzem estrutura espacial fina em praticamente
toda a extensão de GO~e~DF, enquanto o clima permanece quase plano até mesmo na sua própria
escala nativa.

Essa uniformidade é \emph{espacial} (entre lugares, num mesmo instante), não \emph{temporal}
(num mesmo lugar, ao longo do tempo): uma superfície climática pode ser quase plana no espaço e
ainda assim se deslocar substancialmente no tempo, em todo o território de uma vez. É
precisamente esse deslocamento temporal --- não a estrutura espacial fina, que continua a cargo
do terreno e do solo --- que a projeção CMIP6 explora (\S\ref{sec:shift}).

\begin{figure}[H]
\centering
\includegraphics[width=\textwidth]{fig_4_1.png}
\caption{Superfícies de viabilidade do presente (0--1) para os sete segmentos. Os mapas de
classes da FAO para segmentos representativos estão na Figura~\ref{fig:fao-classes} do
Anexo.}
\label{fig:atlas-present}
\end{figure}

\section{Zoneamento agroambiental}\label{sec:zoning-results}

O número de zonas foi selecionado em $k = 2,\dots,10$ pelas estatísticas de silhueta,
Davies--Bouldin e \emph{gap} em conjunto. Neste território fracamente agrupado, a silhueta é
quase plana (0,16--0,18), de modo que nenhum índice isolado é decisivo. Ainda assim, $k = 7$
atinge conjuntamente o melhor Davies--Bouldin (1,56) e a melhor silhueta (empatada), com
estabilidade aceitável em subamostras (índice de Rand ajustado 0,84); o painel completo de
validade e estabilidade é a Figura~\ref{fig:cluster-validity} do Anexo. As sete zonas repartem
o território em parcelas de 19,2\%, 4,8\%, 10,4\%, 10,6\%, 22,1\%, 21,5\% e 11,4\% da área de
interesse (a partir de $n = 14.935$ pontos de treinamento). As médias biofísicas discriminantes
por zona estão na Tabela~\ref{tab:zone-means} do Anexo.

\begin{table}[H]
\centering
\caption{Viabilidade média por segmento, por zona, com o rótulo de melhor segmento
comparativo.}
\label{tab:zone-suit}
\small
\begin{tabular}{@{}lrrrrrrr@{}}
\toprule
Segmento & Z0 & Z1 & Z2 & Z3 & Z4 & Z5 & Z6 \\
\midrule
Soja                  & 0,75 & 0,67 & 0,59 & 0,53 & 0,57 & 0,66 & 0,15 \\
Cana-de-açúcar        & 0,78 & 0,67 & 0,66 & 0,54 & 0,50 & 0,70 & 0,22 \\
Outras culturas       & 0,65 & 0,55 & 0,52 & 0,47 & 0,53 & 0,59 & 0,15 \\
Piscicultura          & 0,23 & 0,14 & 0,15 & 0,35 & 0,18 & 0,12 & 0,02 \\
Pecuária              & 0,56 & 0,56 & 0,66 & 0,54 & 0,49 & 0,53 & 0,51 \\
Conservação           & 0,40 & 0,44 & 0,50 & 0,51 & 0,40 & 0,44 & 0,58 \\
Geração fotovoltaica  & 0,32 & 0,49 & 0,38 & 0,27 & 0,36 & 0,26 & 0,11 \\
\midrule
\textbf{Melhor comparativo} & cana & solar & pecuária & piscicultura & solar & cana & conservação \\
\bottomrule
\end{tabular}
\end{table}

As sete zonas admitem uma interpretação fisiográfica clara, resolvendo cinco vocações
comparativas distintas, incluindo a \textbf{conservação} na Zona~6.

\begin{itemize}
  \item \textbf{Zona 0 (19,2\%) --- planalto de cultivo privilegiado.} Argiloso (36\%), de
    elevação moderada, suavemente inclinado (3,9$^\circ$), com a maior viabilidade de
    culturas de qualquer zona (soja 0,75, cana 0,78); melhor uso comparativo cana-de-açúcar.
    Este é o núcleo agrícola.
  \item \textbf{Zona 1 (4,8\%) --- ombro quente e seco.} Uma zona pequena e a mais seca
    (aridez 0,91), com a maior viabilidade solar (0,49); melhor uso comparativo é a
    \textbf{geração fotovoltaica}.
  \item \textbf{Zona 2 (10,4\%) --- planalto alto e argiloso.} A zona mais alta (862~m) e
    mais argilosa (39\%), mais próxima da água permanente (9,5~km); seu melhor uso comparativo
    é a \textbf{pecuária} (0,66, a mais alta de qualquer zona).
  \item \textbf{Zona 3 (10,6\%) --- baixada ripária.} Densidade de drenagem uma ordem de
    grandeza maior (0,11 contra $\sim$0,01 nas demais) e a maior umidade topográfica (10,9);
    com o melhor uso comparativo sendo a \textbf{piscicultura} (0,35, de longe a maior de qualquer zona).
  \item \textbf{Zona 4 (22,1\%) --- planície baixa, quente e arenosa.} A zona mais baixa
    (370~m), mais plana (3,1$^\circ$), mais quente (26,8~$^\circ$C) e mais arenosa (24\% de
    argila); viabilidade de culturas apenas moderada, estando portanto mais propícia à \textbf{geração
    fotovoltaica} do que ao cultivo. Como mostra a validação (\S\ref{sec:validation-present}), 
      esta é a zona menos produtiva.
  \item \textbf{Zona 5 (21,5\%) --- planalto de cultivo frio e úmido.} Alto (779~m), frio
    (24,3~$^\circ$C), a zona mais úmida (aridez 1,17); a segunda zona de cultivo (soja 0,66,
    cana 0,70), novamente com pendor para a cana.
  \item \textbf{Zona 6 (11,4\%) --- planalto alto e íngreme de conservação.} De longe a mais
    íngreme (13,5$^\circ$, $\sim$três vezes qualquer outra zona), a mais alta em carbono
    orgânico (16,1~g/kg) e em umidade do solo (117~mm). A viabilidade de culturas colapsa
    aqui (soja 0,15), enquanto a viabilidade de conservação é a mais alta de todas (0,58).
    Seu melhor uso comparativo é a \textbf{conservação}, que aqui é também o melhor uso
    \textbf{absoluto} --- a única zona em que a conservação é o segmento mais viável em
    termos absolutos.
\end{itemize}

As zonas são rotuladas pelo melhor segmento comparativo (\S\ref{sec:zoning}), separando o
território em dois planaltos de cultivo (Zonas 0 e 5), um planalto alto de pecuária (Zona 2),
duas planícies quentes com pendor solar (Zonas 1 e 4), uma baixada ripária com pendor para
piscicultura (Zona 3) e um planalto alto de conservação (Zona 6).

\begin{figure}[H]
\centering
\includegraphics[width=0.9\textwidth]{fig_4_3.png}
\caption{Zonas agroambientais ($k = 7$) sobre GO~e~DF.}
\label{fig:zones}
\end{figure}

\section{Mudança projetada pelo clima (QP3)}\label{sec:shift}

A projeção por fatores de mudança (\S\ref{sec:cmip6}) foi concluída para ambos os cenários
(SSP2-4.5, SSP5-8.5) e ambas as janelas (2031--2050, 2051--2070). Para cada combinação, ela
produz uma superfície de viabilidade futura por segmento, a mudança $\Delta S$ em relação à
linha de base do presente rederivada e uma superfície de concordância entre os modelos.
Apenas o bloco climático do conjunto avança: terreno, solo, água, fenologia, acessibilidade
e os fatores laterais estáticos são mantidos fixos. Duas variáveis climáticas (radiação
solar e umidade do solo) também são mantidas na linha de base, o que suprime parte dos
sinais de geração solar e pecuária.

Que apenas o clima avance não contradiz a quase-uniformidade espacial documentada em
\S\ref{sec:atlas-present} (Figura~\ref{fig:scale-mismatch}): hoje, a maioria das pertinências
climáticas está \textbf{saturada} ($\approx$1,0, ex.: a precipitação anual da soja já está bem
dentro da faixa ótima em quase todo o território), e é por isso que o clima quase não
discrimina espaço no presente --- um fator no teto não distingue lugares. Sob os fatores de
mudança do CMIP6, porém, a mesma variável pode sair desse patamar e entrar no trecho inclinado
da curva de pertinência (ou cruzar um limiar), e como a agregação é uma média geométrica
ponderada (produto de pertinências), essa queda se propaga de forma multiplicativa. O
deslocamento climático em si é espacialmente suave (\S\ref{sec:atlas-present}), mas o efeito
resultante em $\Delta S$ \emph{não} é uniforme: um local já perto de uma fronteira de classe
FAO por causa de terreno/solo marginais é empurrado para a classe inferior por um pequeno
recuo climático, enquanto um local com terreno/solo excelentes absorve o mesmo recuo sem mudar
de classe --- daí a realocação de melhor cultura concentrada em fronteiras específicas
(Figura~\ref{fig:best-crop} do Anexo), e não um deslocamento uniforme do mapa inteiro. Em suma:
terreno e solo determinam \emph{onde} o sistema está sensível; o clima fornece o empurrão
\emph{temporal} que decide se essa sensibilidade se materializa.

O sinal projetado é \textbf{modesto em magnitude, mas altamente consistente em direção.} As
mudanças médias territoriais vão de alguns milésimos, nas culturas comerciais resilientes, a
cerca de um décimo, na conservação sob a trajetória mais severa (Tabela~\ref{tab:shift-mean}).
Ainda assim, a fração dos cinco GCMs que concordam quanto ao sinal da mudança situa-se entre
0,92 e 1,00 para todo segmento, cenário e janela sensíveis ao clima. A projeção identifica
\textbf{três segmentos expostos}: a conservação, outras culturas anuais e geração
fotovoltaica; não deixando nenhum segmento estruturalmente invariante. A superfície de
conservação carrega dois fatores climáticos móveis: um termo de viabilidade por aridez e um
termo de calor do trimestre mais quente. Por isso ela responde à projeção e emerge como o
\textbf{segmento mais exposto ao clima de todos}, com a viabilidade da vegetação nativa
declinando sob aquecimento e ressecamento mais acentuadamente que qualquer segmento de
produção.

\begin{table}[H]
\centering
\caption{Mudança média territorial projetada na viabilidade $\Delta S$ (futuro menos
presente rederivado) por segmento. Valores positivos denotam melhora de viabilidade; os
segmentos estão ordenados por exposição.}
\label{tab:shift-mean}
\small
\begin{tabular}{@{}lrrrr@{}}
\toprule
Segmento & SSP2-4.5 2031--50 & SSP2-4.5 2051--70 & SSP5-8.5 2031--50 & SSP5-8.5 2051--70 \\
\midrule
Conservação           & $-0{,}0611$ & $-0{,}0676$ & $-0{,}0729$ & $-0{,}1059$ \\
Outras culturas anuais & $-0{,}0181$ & $-0{,}0316$ & $-0{,}0246$ & $-0{,}0603$ \\
Geração fotovoltaica   & $-0{,}0179$ & $-0{,}0279$ & $-0{,}0234$ & $-0{,}0589$ \\
Soja                   & $+0{,}0006$ & $+0{,}0003$ & $+0{,}0006$ & $-0{,}0033$ \\
Cana-de-açúcar         & $+0{,}0009$ & $+0{,}0011$ & $+0{,}0003$ & $-0{,}0034$ \\
Piscicultura           & $+0{,}0083$ & $+0{,}0088$ & $+0{,}0087$ & $+0{,}0089$ \\
Pecuária               & $-0{,}0000$ & $-0{,}0000$ & $+0{,}0002$ & $-0{,}0010$ \\
\bottomrule
\end{tabular}
\end{table}

Como as pertinências são baseadas em limiares, a medida agronomicamente legível é a mudança
na área classificada como viável ou superior (S1+S2 da FAO, $S \ge 0{,}50$), relatada na
Tabela~\ref{tab:shift-share} como parcela da terra avaliada de cada segmento. Tanto as
parcelas do presente quanto as do futuro são computadas sobre a mesma linha de base rederivada
(presente $=$ futuro $- \Delta S$). Esse presente rederivado coincide com o atlas do presente
exportado até o décimo de ponto percentual, o que confirma a propriedade de cancelamento da
linha de base do mecanismo.

\begin{table}[H]
\centering
\caption{Parcela da terra avaliada de cada segmento classificada como viável ou superior
(S1+S2 da FAO, $S \ge 0{,}50$), presente $\rightarrow$ futuro.}
\label{tab:shift-share}
\small
\begin{tabular}{@{}lrrrrr@{}}
\toprule
Segmento & Presente & SSP2-4.5 2031--50 & SSP2-4.5 2051--70 & SSP5-8.5 2031--50 & SSP5-8.5 2051--70 \\
\midrule
Conservação           & 33,2\% & 18,4\% & 17,1\% & 16,2\% & 11,4\% \\
Outras culturas anuais & 64,2\% & 64,1\% & 63,8\% & 64,0\% & 59,2\% \\
Geração fotovoltaica   & 27,3\% & 23,3\% & 21,2\% & 22,2\% & 17,4\% \\
Soja                   & 62,2\% & 62,2\% & 62,2\% & 62,2\% & 62,1\% \\
Cana-de-açúcar         & 60,8\% & 60,8\% & 60,8\% & 60,8\% & 60,8\% \\
Piscicultura           & 10,0\% & 10,2\% & 10,3\% & 10,3\% & 10,3\% \\
Pecuária               & 63,1\% & 63,1\% & 63,1\% & 63,1\% & 62,7\% \\
\bottomrule
\end{tabular}
\end{table}

A conservação é o segmento mais exposto. Seus fatores de aridez e calor carregam a maior
mudança negativa de qualquer segmento, agravando-se de $-0{,}061$ (SSP2-4.5, curto prazo) para
$-0{,}106$ (SSP5-8.5, fim do século). Sua área viável ou superior mais do que se reduz à
metade, de 33,2\% para 11,4\% sob o cenário e janela mais quente. A vegetação nativa é o uso
da terra menos amortecido pelo clima: os fatores climáticos dos segmentos de produção
situam-se sobre um amplo platô ótimo, pois o manejo os mantém próximos de seu ótimo térmico.
Já os fatores da conservação situam-se sobre a rampa ativa, de modo que o sinal de aquecimento
e ressecamento se transmite diretamente em perda acentuada de terra de alto valor de
conservação.

As outras culturas anuais (milho, algodão, sorgo) são a cultura mais exposta. Sua mudança se
agrava monotonicamente de $-0{,}0181$ para $-0{,}0603$, com a área viável ou superior
contraindo-se de 64,2\% para 59,2\%. Suas pertinências premiam a confiabilidade da chuva e uma
janela térmica limitada, ambas erodidas pela temperatura crescente e pelos graus-dia de
crescimento acumulados. Estes resultados coincidem com a modelagem da pressão exercida pelo
aquecimento e ressecamento sobre a fronteira agrícola do Cerrado, já documentada em escala
regional \citep{Marengo2022}.

A geração fotovoltaica cai de $-0{,}018$ a $-0{,}059$ por sua redução de desempenho por
temperatura de painel, com a área viável ou superior contraindo-se de 27,3\% para 17,4\%: sua
superfície engajada perde área viável com facilidade à medida que as planícies quentes cruzam
o limiar de calor.

Soja e cana-de-açúcar, por sua vez, são resilientes: sua mudança é essencialmente nula, exceto
sob a combinação mais quente, onde se torna levemente negativa ($-0{,}0033$, $-0{,}0034$); sua
área viável ou superior permanece estável até o décimo de ponto percentual em todo o período.
A piscicultura melhora marginalmente (cerca de $+0{,}009$, com a área viável ou superior
passando de 10,0\% para 10,3\%), e a pecuária permanece efetivamente inalterada, dentro de
$\pm 0{,}001$ de zero.

Resolvida por zona (Tabela~\ref{tab:zone-shift} do Anexo), a mudança nas culturas
\textbf{acentua o gradiente agroambiental existente, em vez de deslocá-lo.} As perdas nas
culturas concentram-se na planície baixa, quente e arenosa (Zona~4). Na zona mais quente e
menos produtiva, sob o cenário e janela mais quente, as outras culturas anuais caem $-0{,}127$,
e mesmo a soja --- estável na escala territorial --- declina $-0{,}016$ nessa planície,
a única zona em que cai de forma material.

Nos planaltos de cultivo (Zonas 0 e 5) as mesmas
culturas ficam inalteradas ou melhoram marginalmente, e o planalto alto e íngreme (Zona~6)
mostra apenas pequenos declínios nas culturas. A única perda acentuada da cana-de-açúcar é no
ombro quente e seco (Zona~1, $-0{,}094$), cuja aridez já baixa cruza seu limiar.
O sinal climático, assim, penaliza mais as culturas onde elas já são marginais e
menos onde já são melhores.

Decompondo a mudança em movimentos entre as classes discretas da FAO, confirma-se que a
projeção age quase inteiramente por \emph{rebaixamento} da terra, uma classe de cada vez, e
majoritariamente nos três segmentos expostos (Tabela~\ref{tab:downgrade-share} do Anexo). Para
soja, cana-de-açúcar, piscicultura e pecuária, menos de 1,6\% das células mudam de classe sob
qualquer combinação. Já para a conservação a parcela rebaixada sobe para 43,1\%, para as
outras culturas anuais para 28,0\% e para a geração fotovoltaica para 22,6\% sob a combinação
mais severa. A matriz de transição completa para as outras culturas anuais
(Tabela~\ref{tab:transition-othercrops} do Anexo) mostra que esse rebaixamento se concentra no
topo da escala: sob SSP5-8.5 2051--2070, a classe privilegiada S1 colapsa de 40,7\% para
19,6\% da terra avaliada.

As superfícies de concordância entre os modelos confirmam que essas conclusões direcionais
são robustas ao modelo (concordância média $\geq$ 0,92 para todo segmento). A ordenação de
cenário e janela comporta-se como o raciocínio físico exige. Para todo segmento sensível ao
clima, a mudança é mais pronunciada sob a trajetória de maior emissão e a janela mais tardia:
SSP5-8.5 2051--2070 é uniformemente o caso mais severo, enquanto as duas janelas de curto
prazo diferem pouco. Essa ordenação monótona, junto com a mudança quase nula a partir de
fatores de mudança identidade (\S\ref{sec:cmip6}), é evidência de que o mecanismo responde à
direção e magnitude esperadas.

\begin{figure}[H]
\centering
\includegraphics[width=\textwidth]{fig_4_4.png}
\caption{Mudança projetada na viabilidade $\Delta S$ sob SSP5-8.5, 2051--2070, para os três
segmentos expostos (conservação, outras culturas anuais, geração fotovoltaica) e uma cultura
comercial resiliente (soja). A concordância entre os modelos e a realocação de melhor
cultura estão nas Figuras~\ref{fig:ensemble-agreement} e~\ref{fig:best-crop} do Anexo.}
\label{fig:shift}
\end{figure}

\section{Uso atual da terra e validação (QP2 e validação de QP1)}\label{sec:realized}

A fase final confronta o potencial modelado com o uso atual da terra e com uma medida
independente de produtividade. Três conjuntos de dados do catálogo entram aqui: uma
classificação anual de uso da terra a 30~m agregada à grade, as unidades administrativas de
segundo nível (243 municípios, um dos quais é o Distrito Federal) e os produtos de
produtividade MODIS MOD17 (NPP anual e, para o validador intracultura, GPP da estação de
crescimento).

Cada um é usado \textbf{apenas} como máscara, referência de validação ou área descritiva.
Nenhum entra no conjunto de bandas ou no agrupamento, de modo que a independência do potencial
modelado (QP1) em relação ao uso atual da terra (QP2) seja preservada. Qualquer
concordância entre os dois é, portanto, corroboração do modelo, e não circularidade.

\subsection{Composição do uso atual da terra}\label{sec:realized-composition}

O uso atual da terra de GO~e~DF é dominado por dois usos que juntos ocupam quatro quintos do
território (Tabela~\ref{tab:realized-composition}): pastagem cultivada e em mosaico (45,5\%) e
vegetação nativa (35,0\%). O cultivo anual ocupa cerca de um sexto do estado, majoritariamente
soja (13,8\%), com cana-de-açúcar (2,4\%) e o grupo milho/algodão/arroz (0,6\%) bem menores; a
aquicultura é desprezível em termos de área.

\begin{table}[H]
\centering
\caption{Composição do uso atual da terra de GO~e~DF, como parcela do território.}
\label{tab:realized-composition}
\begin{tabular}{@{}lr@{}}
\toprule
Uso atual & Parcela do território \\
\midrule
Pastagem (cultivada + mosaico agricultura/pastagem)               & 45,5\% \\
Vegetação nativa (floresta, savana, campo, área úmida, afloramento rochoso) & 35,0\% \\
Soja                                                              & 13,8\% \\
Cana-de-açúcar                                                    & 2,4\% \\
Outras culturas anuais (milho, algodão, arroz, mosaico de lavoura) & 0,6\% \\
Não elegível (urbano, mineração, água, não vegetado)              & 2,2\% \\
Outros (silvicultura, culturas perenes)                           & 0,5\% \\
Piscicultura (aquicultura)                                        & $<$0,01\% \\
\bottomrule
\end{tabular}
\end{table}

\subsection{Potencial \emph{versus} realizado (QP2)}\label{sec:pot-vs-real}

\textbf{O zoneamento aprendido separa o uso atual da terra de forma condizente.} 
Caracterizando a composição atual dentro de cada zona biofísica
(Tabela~\ref{tab:zone-realized} do Anexo), o planalto de cultivo privilegiado (Zona~0) carrega
soja em 28,0\% de sua área e cana-de-açúcar em 7,6\% --- as maiores parcelas de culturas de
qualquer zona --- contra a menor cobertura nativa (17,9\%). Já a vegetação nativa sobe para
\textbf{77,7\% do planalto alto e íngreme de conservação (Zona~6)}, justamente a zona cujo
melhor uso modelado é a conservação. A planície quente e arenosa (Zona~4), identificada como
marginal, é dominada por pastagem (60,4\%), com apenas 3,4\% de soja. A validação cruzada mais
forte disponível para a QP1 é justamente esta: uma partição não supervisionada de variáveis
biofísicas reproduz a geografia observada de cultivo, pastejo e vegetação remanescente, e a
zona de conservação é precisamente onde o Cerrado nativo de fato persiste.

\textbf{As culturas atuais ocupam terra de viabilidade substancialmente mais alta que a
média oferecida pelo território.} Onde uma cultura é plantada, a viabilidade média modelada
dessa cultura excede sua média sobre toda a terra convertível: soja 0,929 nos pixels de soja
atual contra 0,774 sobre a terra disponível, cana-de-açúcar 0,948 contra 0,785 e outras
culturas anuais 0,735 contra 0,690. A escolha de sítio revelada concentra-se no extremo alto
da superfície modelada, exatamente como prevê um mapa de potencial válido.

\textbf{Uma folga biofísica ainda separa o potencial do uso atual.} Cruzando as classes da
FAO com o uso atual, obtém-se, por segmento, a parcela do território viável ou superior mas
não atualmente sob esse uso (Tabela~\ref{tab:pot-real-gap}). Para a soja essa parcela
subutilizada é de 48,0\% do território e para a cana-de-açúcar 57,9\%. Como a pastagem é de
longe o maior uso atual, o grosso dessa folga é pastagem existente, e não vegetação nativa.
Isto aponta para um padrão de \textbf{intensificação de terra já convertida}, e não para o
desmatamento.

\begin{table}[H]
\centering
\caption{Lacuna potencial \emph{versus} realizado: parcela do território viável ou superior
(S1+S2 da FAO) para um segmento mas não atualmente sob esse uso.}
\label{tab:pot-real-gap}
\begin{tabular}{@{}lr@{}}
\toprule
Segmento & Parcela subutilizada do território \\
\midrule
Outras culturas anuais & 63,2\% \\
Cana-de-açúcar         & 57,9\% \\
Soja                   & 48,0\% \\
Piscicultura           & 9,9\% \\
\bottomrule
\end{tabular}
\end{table}

\subsection{Validação da viabilidade do presente (QP1)}\label{sec:validation-present}

As superfícies baseadas em conhecimento foram submetidas a três verificações independentes.
Cada uma delas é reforçada com uma validação cruzada espacial, que protege contra a inflação
por autocorrelação espacial a que as métricas de presença estão sujeitas.

\textbf{Discriminação por presença.} Medidas contra a presença atual de culturas, as
superfícies ordenam bem a presença acima do fundo para as duas principais culturas
comerciais (Tabela~\ref{tab:presence-auc}). A soja é validada em ambas as métricas (AUC 0,83;
Boyce 0,66). Sua discriminação é, além disso, \textbf{espacialmente robusta}: recomputando a AUC
dentro de blocos de $\sim$0,5$^\circ$ obtemos um valor médio de 0,79 ($\pm$0,13 entre 50 blocos), 
não estando distante do valor agrupado. A ordenação reflete, portanto, uma capacidade de generalização
e não um resultado dependente da amostra.

A cana-de-açúcar também apresenta uma AUC agrupada forte (0,83) e um Boyce de 0,74, mas sua
AUC média por bloco cai para 0,61 ($\pm$0,13, em apenas seis blocos). Seu alto valor agrupado
é, assim, \textbf{em parte um artefato de ordenamento espacial}: a cana-de-açúcar ocupa poucos
distritos concentrados, e essa é uma ressalva que a métrica agrupada, isolada, teria ocultado.
O grupo das outras culturas, por fim, é validado apenas fracamente (AUC 0,64) --- coerente com
sua pegada atual muito pequena (0,6\%) e sua amostra de presença esparsa.

\begin{table}[H]
\centering
\caption{Discriminação por presença das superfícies de viabilidade contra a presença atual
de culturas (AUC agrupada, AUC média por bloco espacial e índice de Boyce).}
\label{tab:presence-auc}
\begin{tabular}{@{}lrrr@{}}
\toprule
Segmento & AUC (agrupada) & AUC (média por bloco espacial) & Índice de Boyce \\
\midrule
Soja                 & 0,83 & 0,79 $\pm$ 0,13 & 0,66 \\
Cana-de-açúcar       & 0,83 & 0,61 $\pm$ 0,13 & 0,74 \\
Outras culturas anuais & 0,64 & --- & 0,66 \\
\bottomrule
\end{tabular}
\end{table}

\textbf{\emph{Proxy} de produtividade} se refere ao contexto, não ao rendimento. 
Este teste usa a correlação
$\rho$ de Spearman entre a viabilidade média municipal e o \emph{proxy} de NPP anual sobre
terra manejada (Tabela~\ref{tab:spearman-npp} do Anexo), como \emph{contexto} grosseiro de
produtividade, não como um validador de rendimento como já foi explicitado em \S\ref{sec:phaseb}. A
correlação é positiva e significativa para os segmentos de biomassa em pé: pecuária ($+0{,}52$)
e conservação ($+0{,}47$). No caso da conservação, a superfície reconceituada credita a terra
de alta biomassa, acidentada e dominada por vegetação nativa que assimila mais carbono. A
correlação também é positiva para cana-de-açúcar ($+0{,}39$) e solar ($+0{,}27$), mas é nula
para a soja ($+0{,}12$, n.s.) e negativa para a piscicultura ($-0{,}36$).

Essa mistura de sinais tem uma explicação pré-registrada: o NPP anual mede a assimilação de
carbono da vegetação, não o rendimento agronômico, de modo que os segmentos de biomassa em pé
pontuam mais alto. Isso é coerente com a tendência documentada do algoritmo NPP/GPP do MODIS
de estimar a produtividade de forma sistematicamente enviesada por tipo de sítio
\citep{Turner2006}. É coerente também com a literatura sobre a conversão de GPP em rendimento,
em que a correlação entre um \emph{proxy} derivado de GPP e o rendimento reportado varia de
quase unitária na escala de município a marcadamente mais fraca na escala de talhão
\citep{He2018}. É por isso que as métricas de discriminação por presença
(Tabela~\ref{tab:presence-auc}), e não este \emph{proxy}, são os validadores primários das
culturas.

\textbf{Validador de produtividade intracultura.} Este teste avalia a GPP da estação de
crescimento (janela de dossel de outubro a março) \textbf{apenas nos pixels efetivamente
cultivados de cada cultura}, e a correlaciona com a viabilidade dessa cultura \emph{dentro} da
própria cultura --- removendo, assim, as confusões advindas dos ciclos de fenologia e pousio
de coberturas distintas. Para a soja, o gradiente resultante é \textbf{+0,003} (n.s., $n = 641$): 
plano, e não negativo. Isso confirma que uma viabilidade modelada mais alta não está associada a
produtividade \emph{mais baixa}. O gradiente é fraco porque a GPP da estação de crescimento
varia pouco dentro da estação de dossel de uma única cultura. As amostras de cana-de-açúcar e
de outras culturas, por sua vez, são pequenas demais ($n = 119$ e $16$) para serem conclusivas.
As métricas de presença permanecem, portanto, os validadores primários das culturas.

O \emph{proxy} de NPP anual também separa fortemente as zonas. Uma ANOVA de um fator entre as
sete zonas é altamente significativa ($F = 592{,}1$, $p < 0{,}001$): a planície quente e
arenosa (Zona~4) é marcadamente menos produtiva (média do grupo 0,575, em unidades
escalonadas) do que o planalto alto e argiloso (Zona~2, o mais produtivo, com 0,858) e os
planaltos de cultivo (Zonas 0 e 5, com 0,729 e 0,822). O zoneamento, assim, capta um gradiente
real de produtividade realizada, e não apenas um gradiente presente nas variáveis de entrada.

\textbf{Verificação cruzada orientada por dados.} Uma floresta aleatória foi treinada na
presença atual de soja, sobre o mesmo conjunto padronizado de atributos, e comparada célula a
célula com a classificação baseada em conhecimento. Em amostra, a concordância é razoável
($\kappa$ de Cohen $= 0{,}23$): o mapa de conhecimento declara 60,8\% do território como
viável ou superior, contra os 25,4\% a que o classificador atribui alta probabilidade de
presença.

Essa concordância em amostra, porém, é otimista. Sob uma \textbf{validação cruzada de cinco
dobras com blocos espaciais} --- treinando em blocos inteiros de $\sim$0,5$^\circ$ e pontuando
em blocos retidos, de modo que nenhuma célula de teste compartilhe bloco com seus dados de
treinamento --- a concordância fora do bloco cai para $\kappa = 0{,}07$, e a parcela de alta
probabilidade do classificador cai para 6,9\%.

Esse $\kappa$ tão baixo não contradiz a AUC de 0,83 relatada na Tabela~\ref{tab:presence-auc}:
são duas perguntas diferentes sobre o mesmo resultado. A AUC pergunta algo simples --- os pixels
onde a soja é realmente plantada recebem, em geral, uma nota de viabilidade mais alta do que
pixels escolhidos ao acaso? A resposta é sim, de forma consistente em diferentes partes do
estado, por isso ela continua alta. Já o $\kappa$ pergunta algo mais rígido: as duas superfícies
concordam sobre \emph{qual} pedaço do território marcar como viável? Aqui a resposta é não,
porque o mapa de conhecimento marca uma área grande (60,8\%) e a floresta aleatória, sob
validação espacial, marca uma área bem menor (6,9\%). Quando duas áreas comparadas têm tamanhos
tão diferentes, a fórmula do $\kappa$ desconta como ``acerto por sorte'' uma fatia maior da
concordância, e o resultado final cai bastante --- mesmo que o modelo, no fundo, continue
apontando para os lugares certos \citep{Cicchetti1990}. Ou seja, $\kappa = 0{,}07$ não quer dizer
que o modelo erra \emph{onde} a soja tende a ocorrer --- isso já foi respondido pela AUC --- mas
sim que a área que o conhecimento especialista chama de ``viável'' e a área que os dados de fato
ocupam têm tamanhos muito diferentes.

Longe de ser uma falha, essa queda \emph{aguça} a interpretação. Removido o vazamento
espacial, o nicho orientado por dados é muito mais estreito que o envelope biofísico: a maior
parte da terra ``viável'' está genuinamente sem uso. Esse é, precisamente, o sinal potencial
\emph{versus} realizado da QP2, uma vez que terra biofisicamente capaz é deixada ociosa por 
razões de economia, posse da terra ou logística. Como o território é climaticamente quase uniforme, o
solo e o terreno são as únicas alavancas de estreitamento disponíveis, de modo que essa
permissividade residual não pode ser atribuída ao clima. 

Em suma: o mapa baseado em conhecimento responde ``onde a terra é capaz'', e a floresta aleatória 
responde ``onde a soja é de fato plantada''. A divergência entre os dois, controlada espacialmente,
é exatamente a lacuna quantificada em \S\ref{sec:pot-vs-real}.

\subsection{Unidades municipais de decisão}\label{sec:municipal}

Agregando aos 243 municípios obtém-se a visão relevante para a decisão. Por área majoritária,
eles se distribuem por seis das sete zonas: 63 na planície quente e arenosa, 49 no planalto
alto e argiloso, 43 no ombro quente e seco, 32 no planalto de cultivo frio, 30 no planalto de
cultivo privilegiado e 26 na baixada ripária. O planalto alto e íngreme de conservação, por ser
fragmentado e pequeno, nunca é a zona majoritária de um município.

A viabilidade média municipal da soja tem mediana de 0,578 (faixa 0,27--0,90). Trinta e três
municípios alcançam uma média de nível S1 ($\ge 0{,}75$), e 188 dos 243 são, em média, viáveis
ou superiores. Sobrepondo a projeção de meados do século, identificam-se os municípios
vulneráveis: 135 dos 243 enfrentam um declínio projetado das outras culturas mais acentuado
que $-0{,}05$ (média municipal $-0{,}082$, pior caso $-0{,}270$). Esses municípios
concentram-se na planície quente e arenosa, expressando um achado que responde à QP3, indicando que o
nicho de milho--algodão--sorgo se erode mais rápido no terreno mais baixo e já marginal.

Dois \emph{rankings} municipais resumem a leitura relevante para a decisão
(Tabela~\ref{tab:municipal-rank}). Um escore de \textbf{oportunidade} --- viabilidade média da
soja $\times$ a parcela viável mas sem uso --- destaca municípios com alto potencial e ampla
folga em terra já convertida (54 pontuam acima de 0,4). Os líderes combinam viabilidade da
soja perto de 0,77--0,90 com parcelas sem uso acima de 0,68. Um \emph{ranking} de
\textbf{vulnerabilidade} pelo declínio projetado das outras culturas é dominado inteiramente
por municípios de planície quente e arenosa do vale do Araguaia, a noroeste, liderados por
São Miguel do Araguaia ($-0{,}270$).

\begin{table}[H]
\centering
\caption{\emph{Rankings} municipais: cinco primeiros por oportunidade (viabilidade da soja
$\times$ parcela viável mas sem uso) e por vulnerabilidade (mudança projetada $\Delta S$ das
outras culturas, SSP5-8.5 2051--70).}
\label{tab:municipal-rank}
\begin{tabular}{@{}lr@{\hspace{1.5cm}}lr@{}}
\toprule
Oportunidade (município) & Escore & Vulnerabilidade (município) & $\Delta S$ \\
\midrule
Simolândia     & 0,63 & São Miguel do Araguaia & $-0{,}270$ \\
Adelândia      & 0,63 & Novo Planalto           & $-0{,}187$ \\
Gouvelândia    & 0,61 & Nova Crixás             & $-0{,}177$ \\
Firminópolis   & 0,60 & Mundo Novo              & $-0{,}175$ \\
Sanclerlândia  & 0,60 & Bonópolis               & $-0{,}156$ \\
\bottomrule
\end{tabular}
\end{table}

\begin{figure}[H]
\centering
\includegraphics[width=0.85\textwidth]{fig_4_7.png}
\caption{Uso atual da terra de GO~e~DF agregado a 250~m, colorido por função de segmento.}
\label{fig:realized-use}
\end{figure}

\begin{figure}[H]
\centering
\includegraphics[width=0.85\textwidth]{fig_4_9.png}
\caption{Viabilidade média municipal da soja e municípios sinalizados como vulneráveis pelo
declínio projetado das outras culturas. O mapa de subutilização potencial \emph{versus}
realizado da soja está na Figura~\ref{fig:underused-soy} do Anexo.}
\label{fig:municipal}
\end{figure}
